% Full elementary-substitution proof of Kamp III.2, Lemma 9.

\section{Kamp's universal elementary-composition lemma}
\label{sec:annex-p09-kamp-iii2-lemma9}

\begin{annexitem}{annex:p9.kamp.iii2.lemma9}{thm:p09-kamp-iii2-lemma9}{All three families, every labelled disjunct and proper restoration}
\begin{proof}
For each left input choose its supplied finite elementary disjunction
on \((x,z)\). Concatenating all these lists gives one list
\(\rho_1,\ldots,\rho_r\) representing their disjunction \(\Phi\)
under \(x<z\) over \(\mathsf{Com}\). Concatenating the
right lists gives \(\sigma_1,\ldots,\sigma_s\) representing
\(\Psi\) under \(z<y\). When an input family is empty,
its list is empty and its contribution is false. Each nonempty
elementary disjunct has a finite labelled-word presentation by
\cref{thm:p09-kamp-elementary-intervals}.

The finite disjunction of the unary inputs has an elementary history,
because \(\operatorname{E}\) is closed under finite nonempty
disjunction. Its empty case is \(F^K_z\), an allowed At seed.
Thus the formula \(\chi\) used as the single point label
has an actual \(\operatorname{E}(z)\) construction in every case.

Fix a complete linear structure and endpoint assignment. For every
tested \(x<z<y\), both interval guards hold. The supplied
equivalences may therefore be applied at this value of \(z\),
with every original endpoint retained. Consequently the original
universal formula is equivalent over \(\mathsf{Com}\) to
\[
 \forall z\left(x<z<y\to
   \chi(z)\lor\bigvee_{i=1}^r\rho_i(x,z)
                    \lor\bigvee_{j=1}^s\sigma_j(z,y)\right).
 \tag{1}
\]
This argument applies pointwise within the universal guard. It does
not replace a complete-order interval representation by an unguarded
equivalence on another order class.

Choose the finitely many fresh letters prescribed in the body, all
absent from the original unary label and every elementary word label.
The left skeletons belong to \(\operatorname{Bp}(x,z)\), the
right skeletons to \(\operatorname{Bp}(z,y)\), and the one
fresh point atom belongs to \(\operatorname{At}(z)\).
Thus their finite family has exactly the hypotheses of
\cref{thm:p09-kamp-iii2-lemma8}. That theorem constructs a
finite \(\operatorname{Ep}(x,y)\) representation for its universal
formula. This remains applicable when no interval skeleton is present:
the family still contains its point atom, to be replaced by the
appropriate true or false elementary label.

Apply the interval-substitution theorem once for each fresh letter.
Its replacement has exactly one free argument and belongs to
\(\operatorname{E}\). Every new bound variable is chosen fresh
from the enclosing formulas; the outer \(\forall z\) retains
its intended binding of each inserted argument at \(z\), and
the interval patterns retain the analogous bindings at their internal
positions. Each substitution transforms every current elementary
disjunct into an actual elementary disjunct and preserves its
\(\mathsf{Com}\) equivalence under \(x<y\).

All targets are fresh from all replacements. Therefore the chosen
finite sequence restores exactly the original label at each skeleton
occurrence and \(\chi\) at the point atom. No subsequent step
recursively replaces a letter inside an already inserted label.
The resulting universal matrix is exactly (1), up to alphabetic
variation of its internal bound names. This verifies the output
construction rather than assuming that an elementary word is literally
a basic pattern.

Combining the equivalence with (1) proves the claimed property
\(I(x,y)\). If an input interval representation is empty,
its false contribution can either remain an empty list or be written
using the explicit false elementary word; the two choices agree under
the relevant strict guard. If every family is empty, (1) is
\(\forall z\,(x<z<y\to F^K_z)\). It is true precisely
when the tested open interval is empty. The construction preserves
that condition and does not introduce an existential witness.
Every operation is on finite lists or finite substitution histories,
so the construction terminates.

In the four-point example the tested points are exactly 1 and 2.
At 1, the left pattern has \(A(0)\), \(C(1)\) and an
empty \(B\) gap. The right pattern fails because \(D(1)\)
is false. At 2, the left fails because \(C(2)\) is false,
while the right has \(D(2)\), \(F(3)\) and an empty
\(E\) gap. The unary input is false at both points. Hence
the full disjunction holds at every tested point, but neither
individual interval disjunct does. For adjacent outer endpoints
there is no test point at all; in the all-empty-family case the
universal is therefore true. At \((0,3)\), either of the
two interior points makes its false matrix fail. These checks
verify the stated examples and the universal/disjunction distinction
\cite[ch.~III, pp.~92--93]{Kamp1968}.
\end{proof}
\end{annexitem}
